\section{Known defects, invalid runs and untested items}\label{app:honesty}

This appendix collects everything that a reader should know before trusting a number in this report.

\subsection{Invalid or superseded results (kept for the record)}
\begin{enumerate}
\item \textbf{First perception measurement (10 November, 20:31--20:52) invalid.} Stale camera mounts (Section~\ref{sec:perception}); outputs renamed \code{*.SUSPECT_*}; corrected measurement used instead; the dependent \code{hurt} evaluation was discarded and rerun.
\item \textbf{Earlier retention table possibly biased.} It was produced with the same helper; the corrected measurement agrees in order of magnitude for three conditions (0.31, 0.15, 0.014 against 0.27, 0.16, 0.04) but the table was not re-fitted. All weather results of the dropout proxy carry this caveat.
\item \textbf{Bird's-eye view version 1 discarded}: it contained a constant wedge plus detections; one unique grid over 34 frames. Replaced by a camera-derived seven-channel view (version 2), which was in turn \emph{not used} for the final results (the grid of the teammates' repository is). The three behaviour-cloning models on version 2 are kept in \pth{trained_models/02_superseded_early_experiments/}.
\item \textbf{Reward audit.} Version 1 of the reward audit was confounded (the cloned steering drifted at 2\,m/s); version 2 revealed a genuine flaw (a scripted ``fast'' behaviour scored 162 against 144 for the cloned policy); after capping the progress credit and adding an overspeed penalty, version 3 passed (cloned 144.0, crawl 36.0, stationary $-40.0$, fast $-306.6$) on a single seed and a short horizon. This is a sanity check, not a validation.
\item \textbf{Rain-to-friction coupling} in the early environment (declared assumption) affected the first reinforcement-learning attempts and was removed afterwards.
\item \textbf{Colour order.} Frames from MetaDrive are blue-green-red throughout; the fog airlight constant, written for red-green-blue, is therefore slightly warm instead of cool. Cosmetic and consistent between training and evaluation.
\end{enumerate}

\subsection{Items that exist but were not evaluated or finished}
\begin{itemize}
\item \textbf{The student of the collision-aware teacher}: never evaluated on the 200-episode block.
\item \textbf{The unmodified-grid student} and the other input designs: only 10 and 40 episodes; not evaluated on the 200-episode block, so the contribution of the map overlay at that sample size is unknown.
\item \textbf{attempt-4 model with weather input} (reinforcement learning with a weather input): the process ended after step 0 without an error message. No result.
\item \textbf{Image-only model}: training finished, then the program stopped with a serialisation error while writing its final record (\code{Object of type function is not JSON serializable}); the checkpoint \code{best.pt} had been saved before that and is the one evaluated. It was evaluated on 200 clear-weather episodes on 12 November 2026 (Section~\ref{sec:headline}); no weather or corruption experiment exists for it.
\item The no-traffic baseline has the 200-episode clear cell only (the fog cell is partial; rain cells missing).
\item Caution models at 200 episodes in weather: not run (the caution direction was dropped).
\end{itemize}

\subsection{Corrections made when this version of the report was written}
Version 1 of this report (kept unchanged in \pth{archive/report_version_1_engineering_log/}) contained statements that were corrected against the run files:
\begin{itemize}
\item \textbf{Reinforcement-learning attempt 2} was described as fine-tuning the plain-grid model. The run files show that it was the fine-tuning of the proposed model with a weather input (run folder \pth{pilot_cond}); its checkpoint exists. The result (0.675 before and after) is unchanged.
\item \textbf{Attempt 1 training log.} Version 1 quoted a log with returns of $-378$ to $-168$ and no success. The log that belongs to the saved checkpoints shows training success of 0.59 to 0.45 (C2) and 0.62 to 0.60 (Probe C); both logs are now stated (Section~\ref{sec:rl}).
\item \textbf{Attempt 3 learning rate} was stated as $10^{-4}$; it was $5\cdot10^{-5}$ ($10^{-4}$ was attempt 4).
\item \textbf{Departure from a GitHub rule} (exact vehicle positions fed to the learned model through the surrogate) was not stated prominently; it is now stated in Section~\ref{sec:github}.
\end{itemize}

\subsection{Limits of the evidence}
\begin{itemize}
\item \textbf{One training seed per model.} No between-seed variance is known. Differences between models that are small (for example the proposed model against the caution control) may reflect seed variation.
\item \textbf{The policies are evaluated on a surrogate of the camera pipeline}, not on the live pipeline. The equivalence test of the optimised path covered 70 frames of which one had an occupied cell.
\item \textbf{Weather is image corruption only} (and, in the surrogate, detector dropout and range errors); it never changes the vehicle dynamics. The conclusions about weather do not apply to physical rain or fog.
\item \textbf{Only MetaDrive's own rule-based traffic}; one behaviour model. Sparse traffic (density 0.05) is the only regime in which the main results hold.
\item \textbf{Energy} is an estimate from a declared virtual vehicle without regeneration.
\item \textbf{Seeds 5000--5003} were used once by a superseded early pilot; no reported model was selected with them.
\item \textbf{Non-determinism of traffic} across process histories (Section~\ref{sec:repro}).
\item \textbf{The validation rule keeps the earliest of tied rounds}, so several ``dataset-aggregation'' models are in fact round-0 behaviour-cloned models (Section~\ref{sec:catalogue}).
\item \textbf{Literature:} full texts were not read; two statements rest on search summaries and are marked \emph{(unverified)}.
\end{itemize}

\subsection{How the numbers in this report were produced}
Every table and figure is generated by a script from the saved result files; nothing was typed in by hand except the two gap series of Figure~\ref{fig:fidelity} (copied from the diagnosis note, which is archived in \pth{report/data/}) and the perception table (copied from the measurement's own table file). Statistical tests (Wilson intervals, exact McNemar, paired bootstrap) are implemented in \pth{report/make_report_figures.py} and agree with the earlier table \pth{archive/paper_earlier_drafts/tables/table_final_200seed.md} where the two overlap
(for example $p=0.0007$ and $p<0.0001$ for the proposed model against the two baselines; the bootstrap intervals differ in the third decimal because of resampling).
